% !TEX spellcheck = en_US



% ~~~~~~~~~~~~~~~~~~~~~~~~~~~~~~~~~~~~~~~ V
% (NC) 2026 Haluk Bingol
% github.com/halukbingol/LaTeX-Templates
%
% Licensees may copy, distribute, display, and perform the work and make derivative works 
% and remixes based on it only for non-commercial purposes. 
% ~~~~~~~~~~~~~~~~~~~~~~~~~~~~~~~~~~~~~~~ A




% =====================================================================
%  Java on the Command Line -- Freshman Level: Your First Programs in a Terminal
%  ---------------------------------------------------------------------
%  Built on hbCisLaTeXTemplate.tex with the libraries in ../../hblib/.
%  Programs, terminal sessions and their real outputs are in codeoutput/:
%     codeoutput/<Name>.java  a program        codeoutput/<Name>.sh   a session
%     codeoutput/<Name>/      the session's project folder
%     codeoutput/<Name>.txt   what it printed, made by:  bash run_examples.sh
% =====================================================================

\documentclass[aspectratio=169,10pt,compress]{beamer}

% ~~~~~~~~~~~~~~~~~~~~~~~~~~~~~~~~~~~~~~~~~~~~~~~~~~~~~~~~~~~~~~~~~ V hblib
	\usepackage{../../hblib/hblibPresentation} % lib presentation: theme, i/N, \hSection, algorithm2e, ...
	\usepackage{../../hblib/hblibCodeoutput} % lib code + output: \lstcode, \lstterminal, ...
	\usepackage{../../hblib/hblibDatastructures} % lib data structures: \hString, \hStack, \hQueue
% ~~~~~~~~~~~~~~~~~~~~~~~~~~~~~~~~~~~~~~~~~~~~~~~~~~~~~~~~~~~~~~~~~ A hblib


% xcolor, array (and the L{width} column), listings, tikz and algorithm2e come from the hblib libraries
\usepackage[T1]{fontenc}
\usepackage{lmodern}
\usepackage{booktabs}

\definecolor{gitred}{RGB}{240,80,50}
\newcommand{\win}{\textbf{Windows}}
\newcommand{\mac}{\textbf{macOS}}

% ---------------------------------------------------------------------
\title{Java on the Command Line}
\subtitle{Freshman Level: Your First Programs in a Terminal}
\author[Bingol]{Haluk O. Bingol}

\institute[FCIS]{
    Faculty of Computer and Information Sciences\\
    Yeditepe University
}
\date{
	{\tiny \hbTimeStamp}
}

\begin{document}




% ~~~~~~~~~~~~~~~~~~~~~~~~~~~~~~~~~~~~~~ V frm
\begin{frame}[t,fragile]
  \titlepage
\end{frame}
% ~~~~~~~~~~~~~~~~~~~~~~~~~~~~~~~~~~~~~~ A frm




% ~~~~~~~~~~~~~~~~~~~~~~~~~~~~~~~~~~~~~~ V frm
\begin{frame}[t,fragile] \frametitle{Outline}
  \begin{columns}[T]
    \begin{column}{0.48\textwidth}
      \tableofcontents[sections={1-3}, hideallsubsections]
    \end{column}
    \begin{column}{0.48\textwidth}
      \tableofcontents[sections={4-6}, hideallsubsections]
    \end{column}
  \end{columns}
\end{frame}
% ~~~~~~~~~~~~~~~~~~~~~~~~~~~~~~~~~~~~~~ A frm

% ~~~~~~~~~~~~~~~~~~~~~~~~~~~~~~~~~~~~~~ V
\hSection[Install]{Installing Java}{}{}
% ~~~~~~~~~~~~~~~~~~~~~~~~~~~~~~~~~~~~~~ A

% ~~~~~~~~~~~~~~~~~~~~~~~~~~~~~~~~~~~~~~ V
\subsection[Need]{What you need}
% ~~~~~~~~~~~~~~~~~~~~~~~~~~~~~~~~~~~~~~ A




% ~~~~~~~~~~~~~~~~~~~~~~~~~~~~~~~~~~~~~~ V frm
\begin{frame}[t,fragile] \frametitle{What you need}
  \begin{columns}[T]
    \begin{column}{0.5\textwidth}
      \begin{enumerate}
        \item A \textbf{JDK} (Java Development Kit): it contains \texttt{javac} (compiler) and \texttt{java} (runner)
        \item A \textbf{terminal}: \win\ Command Prompt or PowerShell, \mac\ Terminal
        \item A \textbf{text editor} for plain text: VS Code, Notepad++, Sublime, or the built-in editors (with care)
      \end{enumerate}
    \end{column}
    \begin{column}{0.46\textwidth}
      \begin{block}{Which JDK?}
        \small Eclipse \textbf{Temurin 21} (or 25), from \texttt{adoptium.net}. Free, and the installer sets everything up.
      \end{block}
      \begin{alertblock}{JDK, not JRE}
        \small A JRE can only run programs; it has no \texttt{javac}.
      \end{alertblock}
      \begin{block}{Why the terminal if IDEs exist?}
        \small You see exactly what happens; error messages are readable; it works on every computer, server and exam machine.
      \end{block}
    \end{column}
  \end{columns}
\end{frame}
% ~~~~~~~~~~~~~~~~~~~~~~~~~~~~~~~~~~~~~~ A frm

% ~~~~~~~~~~~~~~~~~~~~~~~~~~~~~~~~~~~~~~ V
\subsection[Windows]{Windows}
% ~~~~~~~~~~~~~~~~~~~~~~~~~~~~~~~~~~~~~~ A




% ~~~~~~~~~~~~~~~~~~~~~~~~~~~~~~~~~~~~~~ V frm
\begin{frame}[t,fragile] \frametitle{Installing on Windows}
  \begin{columns}[T]
    \begin{column}{0.52\textwidth}
      \begin{enumerate}\small
        \item Open \texttt{adoptium.net}, choose \textbf{Temurin 21 (LTS)}, \textbf{Windows}, \textbf{x64}, \textbf{JDK}, \texttt{.msi}
        \item Run the \texttt{.msi}. On \emph{Custom Setup} make sure these are enabled:
          \begin{itemize}\footnotesize
            \item \emph{Add to PATH}
            \item \emph{Set JAVA\_HOME variable}
          \end{itemize}
        \item Finish, then open a \textbf{new} Command Prompt (\texttt{Win} key, type \texttt{cmd}, Enter)
        \item Check: \hCode{java -version} and \hCode{javac -version}
      \end{enumerate}
      \footnotesize Or in one command (Windows 10/11): \hCode{winget install EclipseAdoptium.Temurin.21.JDK}
    \end{column}
    \begin{column}{0.44\textwidth}
      \begin{alertblock}{\texttt{'javac' is not recognized \dots}}
        \footnotesize
        \begin{itemize}
          \item Did you open a \textbf{new} window after installing?
          \item \texttt{Add to PATH} missed: \emph{Start}, type \emph{environment variables}, \emph{Edit the system environment variables},
                \emph{Environment Variables\dots}, select \texttt{Path}, \emph{Edit}, \emph{New}:\\
                \texttt{C:\textbackslash Program Files\textbackslash Eclipse Adoptium\textbackslash jdk-21\dots\textbackslash bin}
          \item Check with \hCode{where javac}
        \end{itemize}
      \end{alertblock}
    \end{column}
  \end{columns}
\end{frame}
% ~~~~~~~~~~~~~~~~~~~~~~~~~~~~~~~~~~~~~~ A frm

% ~~~~~~~~~~~~~~~~~~~~~~~~~~~~~~~~~~~~~~ V
\subsection[macOS]{macOS}
% ~~~~~~~~~~~~~~~~~~~~~~~~~~~~~~~~~~~~~~ A




% ~~~~~~~~~~~~~~~~~~~~~~~~~~~~~~~~~~~~~~ V frm
\begin{frame}[t,fragile] \frametitle{Installing on macOS}
  \begin{columns}[T]
    \begin{column}{0.52\textwidth}
      \begin{enumerate}\small
        \item Which chip? \emph{Apple menu, About This Mac}: \textbf{Apple M1/M2/\dots} $\Rightarrow$ \texttt{aarch64}; \textbf{Intel} $\Rightarrow$ \texttt{x64}
        \item Open \texttt{adoptium.net}, choose \textbf{Temurin 21}, \textbf{macOS}, your chip, \textbf{JDK}, \texttt{.pkg}
        \item Run the \texttt{.pkg} and follow the steps
        \item Open \textbf{Terminal} (\texttt{Cmd}+\texttt{Space}, type \emph{Terminal}, Enter)
        \item Check: \hCode{java -version} and \hCode{javac -version}
      \end{enumerate}
      \footnotesize With Homebrew: \hCode{brew install --cask temurin@21}
    \end{column}
    \begin{column}{0.44\textwidth}
      \begin{alertblock}{``Unable to locate a Java Runtime''}
        \footnotesize macOS has a \texttt{java} command that only says Java is missing. Install the JDK, then open a new Terminal window.
      \end{alertblock}
      \begin{block}{Several JDKs?}
        \footnotesize \hCode{/usr/libexec/java\_home -V} lists them. They live in \texttt{/Library/Java/JavaVirtualMachines/}.
      \end{block}
    \end{column}
  \end{columns}
\end{frame}
% ~~~~~~~~~~~~~~~~~~~~~~~~~~~~~~~~~~~~~~ A frm

% ~~~~~~~~~~~~~~~~~~~~~~~~~~~~~~~~~~~~~~ V
\subsection[Check]{Checking}
% ~~~~~~~~~~~~~~~~~~~~~~~~~~~~~~~~~~~~~~ A




% ~~~~~~~~~~~~~~~~~~~~~~~~~~~~~~~~~~~~~~ V frm
\begin{frame}[t,fragile] \frametitle{Checking the installation}
  \small Type both commands; each must print a version, not an error:

  \lstterminal[basicstyle=\ttfamily\scriptsize]{codeoutput/Check.txt}

  \vspace{0.1em}
  \small
  Recorded on Linux with Ubuntu's OpenJDK; with Temurin on Windows or macOS the first line says \texttt{openjdk version "21\dots"} and the
  next lines mention \texttt{Temurin}. The version numbers of \texttt{java} and \texttt{javac} should match.
\end{frame}
% ~~~~~~~~~~~~~~~~~~~~~~~~~~~~~~~~~~~~~~ A frm

% ~~~~~~~~~~~~~~~~~~~~~~~~~~~~~~~~~~~~~~ V
\hSection[Terminal]{Living in the Terminal}{}{}
% ~~~~~~~~~~~~~~~~~~~~~~~~~~~~~~~~~~~~~~ A

% ~~~~~~~~~~~~~~~~~~~~~~~~~~~~~~~~~~~~~~ V
\subsection[Moving]{Moving around}
% ~~~~~~~~~~~~~~~~~~~~~~~~~~~~~~~~~~~~~~ A




% ~~~~~~~~~~~~~~~~~~~~~~~~~~~~~~~~~~~~~~ V frm
\begin{frame}[t,fragile] \frametitle{Moving around in the terminal}
  \begin{columns}[T]
    \begin{column}{0.48\textwidth}
      \begin{block}{Windows (Command Prompt)}


% +++++++++++++++++++++++++++++++++++++++ V lst
%: -Lst.
\begin{lstlisting}[style=codesnippet, language={}, basicstyle=\ttfamily\scriptsize]
C:\Users\ayse> cd Desktop
C:\Users\ayse\Desktop> mkdir cse101
C:\Users\ayse\Desktop> cd cse101
C:\Users\ayse\Desktop\cse101> dir
C:\Users\ayse\Desktop\cse101> cd ..
C:\Users\ayse\Desktop> cd
C:\Users\ayse\Desktop
\end{lstlisting}
% +++++++++++++++++++++++++++++++++++++++ A lst


      \end{block}
    \end{column}
    \begin{column}{0.48\textwidth}
      \begin{block}{macOS (Terminal)}


% +++++++++++++++++++++++++++++++++++++++ V lst
%: -Lst.
\begin{lstlisting}[style=codesnippet, language=bash, basicstyle=\ttfamily\scriptsize]
ayse@mac ~ % cd Desktop
ayse@mac Desktop % mkdir cse101
ayse@mac Desktop % cd cse101
ayse@mac cse101 % ls
ayse@mac cse101 % cd ..
ayse@mac Desktop % pwd
/Users/ayse/Desktop
\end{lstlisting}
% +++++++++++++++++++++++++++++++++++++++ A lst


      \end{block}
    \end{column}
  \end{columns}

  \vspace{0.2em}
  \footnotesize
  \begin{tabular}{@{}lll@{}}
    \toprule
    \textbf{Want to}       & \textbf{Windows}  & \textbf{macOS} \\
    \midrule
    where am I?            & \texttt{cd}       & \texttt{pwd} \\
    list files             & \texttt{dir}      & \texttt{ls} \\
    go into / up           & \texttt{cd name} / \texttt{cd ..} & \texttt{cd name} / \texttt{cd ..} \\
    show a file            & \texttt{type Hello.java} & \texttt{cat Hello.java} \\
    \bottomrule
  \end{tabular}
  \quad
  \begin{minipage}[t]{0.42\textwidth}\footnotesize
    \textbf{Time savers} (both): \texttt{Tab} completes names; $\uparrow$ repeats earlier commands;
    drag a folder onto the window to paste its path.
  \end{minipage}
\end{frame}
% ~~~~~~~~~~~~~~~~~~~~~~~~~~~~~~~~~~~~~~ A frm

% ~~~~~~~~~~~~~~~~~~~~~~~~~~~~~~~~~~~~~~ V
\subsection[Editor]{A plain-text editor}
% ~~~~~~~~~~~~~~~~~~~~~~~~~~~~~~~~~~~~~~ A




% ~~~~~~~~~~~~~~~~~~~~~~~~~~~~~~~~~~~~~~ V frm
\begin{frame}[t,fragile] \frametitle{Writing the file: plain text, right name}
  \begin{columns}[T]
    \begin{column}{0.48\textwidth}
      \begin{block}{Windows Notepad}
        \footnotesize
        \begin{itemize}
          \item \emph{Save as type}: \textbf{All files (*.*)}, name \texttt{Hello.java}; otherwise you get \texttt{Hello.java.txt}
          \item Turn on \emph{File name extensions} in File Explorer (\emph{View}) to see the real name
          \item Encoding: \textbf{UTF-8} (the default in current Windows)
        \end{itemize}
      \end{block}
    \end{column}
    \begin{column}{0.48\textwidth}
      \begin{block}{macOS TextEdit}
        \footnotesize
        \begin{itemize}
          \item \emph{Format, Make Plain Text} (\texttt{Shift}+\texttt{Cmd}+\texttt{T}) before typing; rich text breaks \texttt{javac}
          \item Turn off smart quotes (\emph{Edit, Substitutions}): Java needs straight \texttt{"}, not curly quotes
          \item When saving, untick \emph{If no extension is provided, use .txt}
        \end{itemize}
      \end{block}
    \end{column}
  \end{columns}
  \vspace{0.3em}
  \begin{alertblock}{Easier: a code editor}
    \small VS Code, Notepad++ (Windows) or Sublime Text: plain text by default, colours the code, has a built-in terminal.
  \end{alertblock}
\end{frame}
% ~~~~~~~~~~~~~~~~~~~~~~~~~~~~~~~~~~~~~~ A frm

% ~~~~~~~~~~~~~~~~~~~~~~~~~~~~~~~~~~~~~~ V
\hSection[First]{Your First Program}{}{}
% ~~~~~~~~~~~~~~~~~~~~~~~~~~~~~~~~~~~~~~ A

% ~~~~~~~~~~~~~~~~~~~~~~~~~~~~~~~~~~~~~~ V
\subsection[Compile]{Compile and run}
% ~~~~~~~~~~~~~~~~~~~~~~~~~~~~~~~~~~~~~~ A




% ~~~~~~~~~~~~~~~~~~~~~~~~~~~~~~~~~~~~~~ V frm
\begin{frame}[t,fragile] \frametitle{Compile and run}
  \begin{columns}[T]
    \begin{column}{0.5\textwidth}
      \lstcode[basicstyle=\ttfamily\scriptsize]{codeoutput/First/Hello.java}{Java}
    \end{column}
    \begin{column}{0.48\textwidth}
      \uncover<2->{\lstterminal[basicstyle=\ttfamily\scriptsize]{codeoutput/First.txt}}
    \end{column}
  \end{columns}

  \vspace{0.1em}
  \small \hCode{javac Hello.java} (with \texttt{.java}) creates \texttt{Hello.class}; \hCode{java Hello} (no extension) runs it.
  No output from \texttt{javac} means \textbf{success}.
\end{frame}
% ~~~~~~~~~~~~~~~~~~~~~~~~~~~~~~~~~~~~~~ A frm

% ~~~~~~~~~~~~~~~~~~~~~~~~~~~~~~~~~~~~~~ V
\subsection[Cycle]{Edit, compile, run}
% ~~~~~~~~~~~~~~~~~~~~~~~~~~~~~~~~~~~~~~ A




% ~~~~~~~~~~~~~~~~~~~~~~~~~~~~~~~~~~~~~~ V frm
\begin{frame}[t,fragile] \frametitle{The edit--compile--run cycle}
  \begin{columns}[T]
    \begin{column}{0.42\textwidth}
      \begin{center}
      \begin{tikzpicture}[st/.style={draw, thick, rounded corners, minimum width=2.6cm, minimum height=0.8cm, font=\small, align=center},
                          ar/.style={-{Stealth[length=2.5mm]}, very thick, gray!70}]
        \node[st, fill=blue!8]   (e) at (90:1.6)  {\textbf{edit}\\ \scriptsize save \texttt{Hello.java}};
        \node[st, fill=orange!12](c) at (330:1.6) {\textbf{compile}\\ \scriptsize \texttt{javac Hello.java}};
        \node[st, fill=green!12] (r) at (210:1.6) {\textbf{run}\\ \scriptsize \texttt{java Hello}};
        \draw[ar] (e) to[bend left=25] (c);
        \draw[ar] (c) to[bend left=25] (r);
        \draw[ar] (r) to[bend left=25] (e);
        \draw[ar, red!60] (c.north east) to[out=60, in=0] node[right, font=\scriptsize, text=red!70!black]{errors} (e.east);
      \end{tikzpicture}
      \end{center}
    \end{column}
    \begin{column}{0.56\textwidth}
      \uncover<2->{\lstterminal[basicstyle=\ttfamily\tiny]{codeoutput/Recompile.txt}}
      \uncover<2->{\footnotesize Between the first and second \texttt{java Hello} the file was changed in the editor and saved,
      but \textbf{not compiled}: \texttt{java} still ran the old \texttt{Hello.class}.}
    \end{column}
  \end{columns}
\end{frame}
% ~~~~~~~~~~~~~~~~~~~~~~~~~~~~~~~~~~~~~~ A frm

% ~~~~~~~~~~~~~~~~~~~~~~~~~~~~~~~~~~~~~~ V
\subsection[Input]{Input and arguments}
% ~~~~~~~~~~~~~~~~~~~~~~~~~~~~~~~~~~~~~~ A




% ~~~~~~~~~~~~~~~~~~~~~~~~~~~~~~~~~~~~~~ V frm
\begin{frame}[t,fragile] \frametitle{Reading input with \texttt{Scanner}}
  \begin{columns}[T]
    \begin{column}{0.56\textwidth}
      \lstcode[basicstyle=\ttfamily\tiny]{codeoutput/Greet/Greet.java}{Java}
    \end{column}
    \begin{column}{0.42\textwidth}
      \uncover<2->{\lstterminal[basicstyle=\ttfamily\tiny]{codeoutput/Greet.txt}}
    \end{column}
  \end{columns}

  \vspace{0.1em}
  \small When \textbf{you} run it, you type the answers after the prompts. Here they came from \texttt{printf \dots\ |} and from a file
  (\texttt{< answers.txt}), so they are not shown on the screen.
\end{frame}
% ~~~~~~~~~~~~~~~~~~~~~~~~~~~~~~~~~~~~~~ A frm




% ~~~~~~~~~~~~~~~~~~~~~~~~~~~~~~~~~~~~~~ V frm
\begin{frame}[t,fragile] \frametitle{Command-line arguments}
  \begin{columns}[T]
    \begin{column}{0.54\textwidth}
      \lstcode[basicstyle=\ttfamily\tiny]{codeoutput/Args/Area.java}{Java}
    \end{column}
    \begin{column}{0.44\textwidth}
      \uncover<2->{\lstterminal[basicstyle=\ttfamily\tiny]{codeoutput/Args.txt}}
    \end{column}
  \end{columns}

  \vspace{0.1em}
  \small Arguments are \textbf{strings}; \texttt{Double.parseDouble} converts them and throws \texttt{NumberFormatException} on bad input.
  Check \texttt{args.length} first and catch the exception.
\end{frame}
% ~~~~~~~~~~~~~~~~~~~~~~~~~~~~~~~~~~~~~~ A frm

% ~~~~~~~~~~~~~~~~~~~~~~~~~~~~~~~~~~~~~~ V
\hSection[Errors]{First-Week Errors}{}{}
% ~~~~~~~~~~~~~~~~~~~~~~~~~~~~~~~~~~~~~~ A

% ~~~~~~~~~~~~~~~~~~~~~~~~~~~~~~~~~~~~~~ V
\subsection[Not found]{Command not found}
% ~~~~~~~~~~~~~~~~~~~~~~~~~~~~~~~~~~~~~~ A




% ~~~~~~~~~~~~~~~~~~~~~~~~~~~~~~~~~~~~~~ V frm
\begin{frame}[t,fragile] \frametitle{Error 0: the command is not found}
  \begin{columns}[T]
    \begin{column}{0.48\textwidth}
      \begin{block}{Windows (Command Prompt)}


% +++++++++++++++++++++++++++++++++++++++ V lst
%: -Lst.
\begin{lstlisting}[style=codesnippet, language={}, basicstyle=\ttfamily\scriptsize]
C:\Users\ayse\cse101> javac Hello.java
'javac' is not recognized as an internal or
external command, operable program or batch file.
\end{lstlisting}
% +++++++++++++++++++++++++++++++++++++++ A lst


      \end{block}
    \end{column}
    \begin{column}{0.48\textwidth}
      \begin{block}{macOS (Terminal)}


% +++++++++++++++++++++++++++++++++++++++ V lst
%: -Lst.
\begin{lstlisting}[style=codesnippet, language=bash, basicstyle=\ttfamily\scriptsize]
ayse@mac cse101 % javac Hello.java
zsh: command not found: javac
\end{lstlisting}
% +++++++++++++++++++++++++++++++++++++++ A lst


      \end{block}
    \end{column}
  \end{columns}

  \vspace{0.3em}
  \begin{itemize}\small
    \item Not a Java error: the \textbf{terminal} cannot find the program
    \item Fix: install the JDK, open a \textbf{new} terminal window, check \texttt{PATH} (Windows: \texttt{where javac}, macOS: \texttt{which javac})
    \item \texttt{java} works but \texttt{javac} does not? A JRE or an old Java is first on the \texttt{PATH}
  \end{itemize}
\end{frame}
% ~~~~~~~~~~~~~~~~~~~~~~~~~~~~~~~~~~~~~~ A frm

% ~~~~~~~~~~~~~~~~~~~~~~~~~~~~~~~~~~~~~~ V
\subsection[Compile errors]{Compile errors}
% ~~~~~~~~~~~~~~~~~~~~~~~~~~~~~~~~~~~~~~ A




% ~~~~~~~~~~~~~~~~~~~~~~~~~~~~~~~~~~~~~~ V frm
\begin{frame}[t,fragile] \frametitle{Compile errors: file name and class name}
  \begin{columns}[T]
    \begin{column}{0.48\textwidth}
      \lstterminal[basicstyle=\ttfamily\scriptsize]{codeoutput/WrongFile.txt}
      \footnotesize A \texttt{public class Hello} must be in \texttt{Hello.java}, with the same capital letters.
    \end{column}
    \begin{column}{0.48\textwidth}
      \lstterminal[basicstyle=\ttfamily\scriptsize]{codeoutput/NoFile.txt}
      \footnotesize A typo in the file name, or the file is in another folder: check with \texttt{ls} (\win: \texttt{dir}).
    \end{column}
  \end{columns}
\end{frame}
% ~~~~~~~~~~~~~~~~~~~~~~~~~~~~~~~~~~~~~~ A frm




% ~~~~~~~~~~~~~~~~~~~~~~~~~~~~~~~~~~~~~~ V frm
\begin{frame}[t,fragile] \frametitle{Compile errors: typos}
  \begin{columns}[T]
    \begin{column}{0.44\textwidth}
      \lstcode[basicstyle=\ttfamily\tiny]{codeoutput/Typos/Hello.java}{Java}
    \end{column}
    \begin{column}{0.54\textwidth}
      \uncover<2->{\lstterminal[basicstyle=\ttfamily\tiny]{codeoutput/Typos.txt}}
    \end{column}
  \end{columns}

  \vspace{0.1em}
  \small Java is \textbf{case-sensitive}: \texttt{system} is not \texttt{System} (``package system does not exist'').
  \textbf{cannot find symbol} means ``I don't know this name'': here a typo, \texttt{printn}. Each error gives file, line and a \texttt{\^{}}.
\end{frame}
% ~~~~~~~~~~~~~~~~~~~~~~~~~~~~~~~~~~~~~~ A frm

% ~~~~~~~~~~~~~~~~~~~~~~~~~~~~~~~~~~~~~~ V
\subsection[Run errors]{Run errors}
% ~~~~~~~~~~~~~~~~~~~~~~~~~~~~~~~~~~~~~~ A




% ~~~~~~~~~~~~~~~~~~~~~~~~~~~~~~~~~~~~~~ V frm
\begin{frame}[t,fragile] \frametitle{Run errors: what you give to \texttt{java}}
  \lstterminal[basicstyle=\ttfamily\scriptsize]{codeoutput/RunErrors.txt}

  \vspace{0.1em}
  \small
  \texttt{java} wants the \textbf{class name}: no \texttt{.class}, same capitals, no typos. On Windows and macOS (case-insensitive disks)
  \texttt{java hello} may instead report \texttt{wrong name: Hello}; the fix is the same.
\end{frame}
% ~~~~~~~~~~~~~~~~~~~~~~~~~~~~~~~~~~~~~~ A frm




% ~~~~~~~~~~~~~~~~~~~~~~~~~~~~~~~~~~~~~~ V frm
\begin{frame}[t,fragile] \frametitle{Run errors: no \texttt{main} method}
  \begin{columns}[T]
    \begin{column}{0.4\textwidth}
      \lstcode[basicstyle=\ttfamily\scriptsize]{codeoutput/NoMain/Tools.java}{Java}
    \end{column}
    \begin{column}{0.58\textwidth}
      \uncover<2->{\lstterminal[basicstyle=\ttfamily\tiny]{codeoutput/NoMain.txt}}
    \end{column}
  \end{columns}

  \vspace{0.1em}
  \small A class can compile without \texttt{main}, but only a class with
  \texttt{public static void main(String[] args)} can be started with \texttt{java}.
\end{frame}
% ~~~~~~~~~~~~~~~~~~~~~~~~~~~~~~~~~~~~~~ A frm

% ~~~~~~~~~~~~~~~~~~~~~~~~~~~~~~~~~~~~~~ V
\hSection[Jar]{Running a Program You Were Given}{}{}
% ~~~~~~~~~~~~~~~~~~~~~~~~~~~~~~~~~~~~~~ A

% ~~~~~~~~~~~~~~~~~~~~~~~~~~~~~~~~~~~~~~ V
\subsection[Run jar]{Running a JAR}
% ~~~~~~~~~~~~~~~~~~~~~~~~~~~~~~~~~~~~~~ A




% ~~~~~~~~~~~~~~~~~~~~~~~~~~~~~~~~~~~~~~ V frm
\begin{frame}[t,fragile] \frametitle{Running a JAR file}
  \small Your instructor gives you \texttt{quiz.jar}. You do not need its source code or \texttt{javac}:

  \lstterminal[basicstyle=\ttfamily\scriptsize]{codeoutput/RunJar.txt}

  \vspace{0.1em}
  \small A \textbf{JAR} packs compiled classes into one file. Double-clicking it may start it without a terminal window,
  so you would not see what it prints: run console programs with \hCode{java -jar} in a terminal (Windows and macOS alike).
\end{frame}
% ~~~~~~~~~~~~~~~~~~~~~~~~~~~~~~~~~~~~~~ A frm

% ~~~~~~~~~~~~~~~~~~~~~~~~~~~~~~~~~~~~~~ V
\hSection[Summary]{Summary}{}{}
% ~~~~~~~~~~~~~~~~~~~~~~~~~~~~~~~~~~~~~~ A

% ~~~~~~~~~~~~~~~~~~~~~~~~~~~~~~~~~~~~~~ V
\subsection[Cheat sheet]{Cheat sheet}
% ~~~~~~~~~~~~~~~~~~~~~~~~~~~~~~~~~~~~~~ A




% ~~~~~~~~~~~~~~~~~~~~~~~~~~~~~~~~~~~~~~ V frm
\begin{frame}[t,fragile] \frametitle{Cheat sheet}
  \footnotesize
  \begin{tabular}{@{}L{0.3\textwidth}L{0.31\textwidth}L{0.31\textwidth}@{}}
    \toprule
    \textbf{Task}                & \textbf{Windows (Command Prompt)} & \textbf{macOS (Terminal)} \\
    \midrule
    check Java                   & \texttt{java -version}            & \texttt{java -version} \\
    go to your folder            & \texttt{cd Desktop\textbackslash cse101} & \texttt{cd Desktop/cse101} \\
    list files                   & \texttt{dir}                      & \texttt{ls} \\
    compile                      & \texttt{javac Hello.java}         & \texttt{javac Hello.java} \\
    run                          & \texttt{java Hello}               & \texttt{java Hello} \\
    run with arguments           & \texttt{java Area 3 4}            & \texttt{java Area 3 4} \\
    input from a file            & \texttt{java Greet < answers.txt} & \texttt{java Greet < answers.txt} \\
    output to a file             & \texttt{java Hello > out.txt}     & \texttt{java Hello > out.txt} \\
    run a JAR                    & \texttt{java -jar quiz.jar}       & \texttt{java -jar quiz.jar} \\
    stop a running program       & \texttt{Ctrl}+\texttt{C}          & \texttt{Ctrl}+\texttt{C} \\
    \bottomrule
  \end{tabular}
\end{frame}
% ~~~~~~~~~~~~~~~~~~~~~~~~~~~~~~~~~~~~~~ A frm

% ~~~~~~~~~~~~~~~~~~~~~~~~~~~~~~~~~~~~~~ V
\subsection[Rules]{Do and do not}
% ~~~~~~~~~~~~~~~~~~~~~~~~~~~~~~~~~~~~~~ A




% ~~~~~~~~~~~~~~~~~~~~~~~~~~~~~~~~~~~~~~ V frm
\begin{frame}[t,fragile] \frametitle{Do and do not}
  \begin{columns}[T]
    \begin{column}{0.48\textwidth}
      \begin{block}{Do}
        \small
        \begin{itemize}
          \item Name the file exactly like the public class
          \item Compile again after every change
          \item Read the \textbf{first} error message, line number included
          \item Keep each exercise in its own folder
          \item Use a plain-text code editor
        \end{itemize}
      \end{block}
    \end{column}
    \begin{column}{0.48\textwidth}
      \begin{alertblock}{Do not}
        \small
        \begin{itemize}
          \item \texttt{java Hello.class} (give the class name: \texttt{java Hello})
          \item \texttt{javac Hello} (it needs \texttt{.java})
          \item Type \texttt{system} or \texttt{Main} where Java expects \texttt{System} or \texttt{main}
          \item Save as \texttt{.txt} or rich text
          \item Work in the Downloads folder
        \end{itemize}
      \end{alertblock}
    \end{column}
  \end{columns}
\end{frame}
% ~~~~~~~~~~~~~~~~~~~~~~~~~~~~~~~~~~~~~~ A frm

% ~~~~~~~~~~~~~~~~~~~~~~~~~~~~~~~~~~~~~~ V
\subsection[Exercises]{Exercises}
% ~~~~~~~~~~~~~~~~~~~~~~~~~~~~~~~~~~~~~~ A




% ~~~~~~~~~~~~~~~~~~~~~~~~~~~~~~~~~~~~~~ V frm
\begin{frame}[t,fragile] \frametitle{Exercises}
  \begin{enumerate}\small
    \item Install the JDK and show \texttt{java -version} and \texttt{javac -version} in a terminal.
    \item Create a folder \texttt{cse101/lab1} from the terminal; write, compile and run \texttt{Hello.java} there.
    \item Make each of the errors on the ``First-Week Errors'' slides on purpose and read the message.
    \item Write \texttt{Rectangle.java} that takes width and height as arguments and prints area and perimeter.
    \item Write \texttt{Average.java} that reads numbers with \texttt{Scanner} until the input ends; test it with \texttt{<} and a file.
    \item Save the output of a program to \texttt{out.txt} with \texttt{>} and open it in your editor.
    \item Run \texttt{quiz.jar} with and without an argument.
  \end{enumerate}
\end{frame}
% ~~~~~~~~~~~~~~~~~~~~~~~~~~~~~~~~~~~~~~ A frm




% ~~~~~~~~~~~~~~~~~~~~~~~~~~~~~~~~~~~~~~ V frm
\begin{frame}[t,fragile]
	\begin{center}
		\vspace{4cm}
		\hRemark{\Huge Questions?}
	\end{center}
\end{frame}
% ~~~~~~~~~~~~~~~~~~~~~~~~~~~~~~~~~~~~~~ A frm



\end{document}
